\section{Method}

% Generated by political_classifier/scripts/07_build_attention_outputs.py.
\input{output/tables/attention/political_corruption_manuscript_values}

\subsection{Data Collection and Political-Corruption Classifier}
\label{sec:method-political-classifier}

News articles were collected as part of the RESPOND multilingual news corpus.
The data-collection workflow queried the News API using country-specific
corruption keyword dictionaries with declensions and morphological variants.
The collection code and keyword-generation notebooks are maintained in the
related RESPOND media repository.\footnote{\url{https://github.com/annekroon/RESPOND_media/tree/main/data-collection/news-collection/news-api}}
The present repository starts from \PCRawQueryN{} exported query results from
Bulgaria, France, Hungary, Italy, the Netherlands, Serbia, Sweden, Ukraine, and
the United Kingdom.

The first stage cleaned and deduplicated each country-level file independently. The combined
title and body field was standardized as \texttt{article\_text}; excess
whitespace and common character-encoding artefacts were normalized. Rows marked
as duplicates by the source data were removed, as were articles with missing
text or fewer than 80 words. Duplicate URIs and exact duplicate normalized texts were then
removed within country using SHA-256 hashes. Prefix-based near-duplicate removal
was not used in the production run because common article leads can occur in
substantively different stories. Country-level processing prevents an article
published in more than one national corpus from being silently assigned to the
first country encountered. The resulting cleaned corruption-query corpus
contained \PCCleanedQueryN{} articles.

\subsubsection{Source Exclusion Screen}

To improve cross-country comparability, the cleaned corpus was screened to
remove outlets explicitly identified as non-eligible. Each country--source pair was evaluated
against three criteria: the outlet had to be genuinely based in the country to
which it was assigned; it had to produce original reporting or exercise
identifiable editorial curation; and it could not function primarily as an
official repository, press-release distributor, or automated aggregator.
Descriptive outlet characteristics, including source type, national versus
regional scope, and broadcast versus print/digital format, were recorded
separately and did not by themselves determine eligibility.

The reviewed source workbook records the final exclusion decision in the
\texttt{conventional\_journalism} field. Country--source pairs explicitly coded
\texttt{No} were removed. Sources coded \texttt{Yes}, left unresolved, or absent
from the workbook were retained; the latter two groups were separately audited.
This exclusion screen reduced the cleaned corruption-query corpus from
\PCCleanedQueryN{} to \PCSourceFilteredN{} articles. The source screen was applied before drawing
the silver-labelled training sample and before scoring the corpus.

\subsubsection{Classifier Development and Validation}

The second stage distinguished articles primarily about political corruption
from articles that mentioned corruption only incidentally or concerned
non-political misconduct. Political corruption was defined as the misuse of
public or political office, authority, or resources for private or political
gain. The positive class included reporting centered on alleged, investigated,
charged, or established corruption involving political officials or public
authority. Articles centered on private fraud, ordinary crime, or misconduct by
non-political actors were coded as negative.

Because manually labelled training data were limited, a country-balanced
candidate set of 500 source-filtered articles per country was labelled through
the UvA LLM proxy using GPT-5.1 with temperature set to 0. After unusable cases
and duplicate article identifiers were removed, the combined silver-labelled
training set contained \PCSilverTrainN{} articles. Before training, URI and
normalized-text hashes were compared with the human benchmark; any overlap was
removed from the candidate set and a zero-overlap assertion was required to
pass. These model-generated labels were used only for training. They were not
treated as human annotations.

Classifier selection used a separate manually annotated benchmark. The
original benchmark contained 452 articles. Because the initial United Kingdom
subset contained only five positive cases, 50 additional UK articles were
manually reviewed and added as a validation-only supplement. The final
benchmark was restricted to the same source-eligible population as the
production corpus, leaving \PCHumanBenchmarkN{} articles, of which
\PCHumanPositiveN{} were political corruption. Neither the original benchmark
nor the UK supplement was used to train the selected silver-label model.

Before model fitting, the eligible benchmark was divided once, using a fixed
random seed and stratification by country and class, into a threshold-calibration
partition (\PCCalibrationN{} articles) and a held-out reporting partition
(\PCHeldOutN{} articles, including \PCHeldOutPositiveN{} positive cases). All
candidate thresholds were selected on the calibration partition. The held-out
partition was opened only to report the performance of the pre-specified final
embedding model and selected threshold. Human-label five-fold cross-validation
models are retained as diagnostics and are explicitly distinguished from the
held-out silver-model estimate in the tables.

The selected classifier combined
\texttt{intfloat/}\allowbreak\texttt{multilingual-e5-large} sentence embeddings with
class-balanced logistic regression. A validation-set threshold sweep selected
a decision threshold of \PCThreshold{} on the calibration partition. On the
held-out reporting partition, the classifier achieved accuracy =
\PCTestAccuracy{}, political-corruption precision = \PCTestPrecision{},
political-corruption recall = \PCTestRecall{}, political-corruption F1 =
\PCTestFOne{}, macro F1 = \PCTestMacroFOne{}, and weighted F1 =
\PCTestWeightedFOne{}. Table~\ref{tab:pc-classifier-comparison-main}
reports the selected classifier's held-out performance. Appendix
Tables~\ref{tab:pc-classifier-comparison-appendix},
\ref{tab:pc-classifier-threshold-sweep},
and~\ref{tab:pc-classifier-country-validation} report the full diagnostic
comparison, threshold sweep, and country-level results, respectively.

\begin{table}
\caption{Validation performance of political-corruption classification models.}
\label{tab:pc-classifier-comparison-main}
\resizebox{\textwidth}{!}{%
\begin{tabular}{llrrrrrrrr}
\toprule
Model & Labels & N train & Thr. & Acc. & PC Prec. & PC Rec. & PC F1 & Macro F1 & Wtd. F1 \\
\midrule
E5-large & Human CV & 502 & 0.500 & 0.817 & 0.634 & 0.823 & 0.716 & 0.790 & 0.823 \\
E5-large & Silver 1+2 & 3982 & 0.400 & 0.863 & 0.731 & 0.809 & 0.768 & 0.835 & 0.865 \\
TF-IDF char. n-grams & Human CV & 502 & 0.500 & 0.777 & 0.601 & 0.610 & 0.606 & 0.725 & 0.777 \\
\bottomrule
\end{tabular}
}
\par\smallskip\footnotesize{Note. PC = political corruption. Precision, recall, and PC F1 refer to the political-corruption class.}
\end{table}


The selected model was applied to all \PCSourceFilteredN{} source-filtered
corruption-query articles. Articles with predicted probabilities at or above
\PCThreshold{} were retained, producing a final political-corruption corpus of
\PCFinalN{} articles, or \PCFinalShare\% of the source-filtered query corpus. These articles form
the input for the article-level content coding and attention analyses described
below.

The attention analysis uses a second, independent News API route. Weekly
country-level counts of all collected news coverage, comprising \PCTotalNewsN{}
news items, provide the denominator for relative attention. The numerator is
the number of source-eligible articles classified as political corruption in
the same country-period. Because the count endpoint does not provide
source-specific records, this total-news denominator cannot be screened with
the outlet-inclusion workbook and represents the broader collected news
universe. Figure~\ref{fig:pc-data-pipeline} distinguishes the article-retrieval
and total-news-count routes. Country-specific corpus-construction counts are
reported in Table~\ref{tab:pc-corpus-construction-country-summary}.

\input{output/tables/attention/figure_political_corruption_data_pipeline_tikz}

\input{output/tables/attention/table_attention_corpus_construction_country_summary}

\subsection{Exploratory Topic Modeling}
\label{sec:method-bertopic}

BERTopic was used as a separate inductive step to inspect recurring substantive
patterns in the final political-corruption corpus. The topic model was not used
to assign the article-level variables, to estimate their validity, or as an
inferential predictor. Its purpose was exploratory: representative articles,
keywords, country distributions, and temporal patterns were examined to
sensitize the later codebook to recurring forms of corruption coverage.

To reduce domination by large countries and years, the topic-model sample was
stratified by country and publication year, with up to 200 political-corruption
articles drawn per non-empty stratum and sampling weights retained. The
multilingual BERTopic specification used
\texttt{intfloat/}\allowbreak\texttt{multilingual-e5-large} embeddings, a minimum
topic size of 10, and automatic topic reduction. GPT-5.1 proposed readable
labels for the discovered clusters. The researchers then defined six broad
interpretive groupings, and GPT-5.1 assigned the fine-grained clusters to those
predefined groupings using the topic keywords and representative documents.
These groupings are analyst-defined summaries of an inductive topic solution,
not additional inductively estimated topics. Prompts, selected examples, raw
responses, model metadata, and run
manifests were retained for audit. The resulting topic solution is reported as
exploratory appendix material and informed, but did not determine, the
substantive coding categories.

\subsection{Article-Level Codebook Development}
\label{sec:method-codebook-development}

The unit of analysis for the substantive variables is the news article. The
codebook combined the theoretical distinctions developed above, close reading
of political-corruption articles, patterns encountered in the exploratory topic
model, and iterative review of disagreements between researcher and model
coding.

A country-balanced development sample was drawn from articles already
classified as political corruption. It contained 108 articles, with 12 from
each country. Within country, sampling was spread across publication years and
three political-corruption probability bands so that development work included
both high-confidence cases and cases closer to the \PCThreshold{} corpus-inclusion
threshold. The articles were translated into English with GPT-5.1 for
researcher coding; the annotation interface retained and displayed the original
text alongside the translation so that ambiguous translations could be
checked.

The development sample was used to clarify category boundaries and identify
systematic errors. In particular, review focused on inferred rather than stated
victims, harm caused by investigation or scandal rather than by corruption,
the distinction between individualized cases and systemic explanations, the
location of transnational cases, and overly broad attribution of an
individual's conduct to an organization. Each revision was implemented as a
new prompt version, and earlier model outputs were retained as pilot results
rather than overwritten.

Because these 108 articles informed the codebook and prompts, they are
development data and cannot provide an unbiased final performance estimate.
After the prompt definitions are frozen, final measurement performance will be
evaluated on a separate country-year-stratified human-coded sample from the
\PCFinalN{}-article corpus. Human--model agreement, variable-specific confusion
matrices, and country-level results will be reported on that held-out sample.
Any disagreements used to modify a definition or prompt will remain excluded
from the final validation estimate.

\subsection{Prompt-Based Article Coding}
\label{sec:method-llm-coding}

The four substantive concepts were implemented as separate prompt-based LLM
coding tasks through the UvA LLM proxy. GPT-5.1 was used during iterative
codebook development and for the reported pilot classifications. Alternative
models are considered only after access through the same proxy and an
identical-prompt comparison have been verified; unavailable models are not
treated as evaluated candidates. Before final corpus coding, the selected
model and prompt versions are frozen and recorded in the run manifests, and
performance is assessed on the separate held-out human-coded sample. The
development articles are not used as the final validation sample. Separating
the four tasks prevented one concept's
decision from mechanically determining another. Each request supplied the
publication country and up to the first 6,000 normalized characters of article
text. The prompts instructed the model to use only the article, to code the
article's representation of allegations rather than their truth, and first to
isolate the corruption case from unrelated crimes, controversies, or harms.
Temperature was set to 0 where supported.

Each task required structured JSON containing the categorical label, supporting
textual evidence, a brief rationale, and a confidence score. The production
scripts save the model name, prompt version, prompt, raw response, parsed
response, timestamp, and any error in per-article JSONL audit files. Output is
checkpointed and resumable. The current versioned specifications are
\nolinkurl{victim_visibility_zero_shot_v8},
\nolinkurl{corruption_frame_zero_shot_v3},
\nolinkurl{abroad_case_zero_shot_v3}, and
\nolinkurl{accused_actor_zero_shot_v6}. Earlier outputs are treated as
development pilots and must not be mixed with results produced by these
versions.

\subsection{Article-Level Variables}
\label{sec:method-variables}

\subsubsection{Victim Visibility}

Victim visibility captures who or what the article explicitly represents as
having suffered harm because of the corruption being discussed. A positive
code requires textual support for three elements: an identifiable person,
group, organization, institution, or public interest; a realized or alleged
realized loss, injury, deprivation, or adverse treatment; and a direct
connection between that harm and the corruption. The procedure does not infer a
victim from the offense type, the involvement of public money, or the existence
of a beneficiary.

The variable distinguishes \nolinkurl{no_victim},
\nolinkurl{concrete_victim},
\texttt{institutional\_}\allowbreak\texttt{societal\_victim}, and
\nolinkurl{unclear}.
\nolinkurl{concrete_victim} applies when a bounded person, group, community,
company, association, or similar entity is explicitly described as losing
money, property, rights, services, employment, contracts, or opportunities, or
as being subjected to corruption-related coercion or extortion.
\texttt{institutional\_}\allowbreak\texttt{societal\_victim} applies when
corruption is explicitly represented as harming a public institution, public finances, public service,
democracy, electoral legitimacy, the rule of law, public trust, state capacity,
society, the economy, or development. If both types are explicit, the final
single-label variable gives priority to \nolinkurl{concrete_victim}.
\nolinkurl{unclear} is reserved for genuinely incomplete, contradictory, or
translation-problematic evidence.

The victim task applies a strict causal-source rule. Harm caused by an
investigation, prosecution, trial, resignation, scandal, political controversy,
or institutional response does not count as harm caused by corruption unless
the article separately makes that connection. Possible, intended, or
hypothetical future harm is insufficient. A communicated extortionate demand
or corrupt threat is an exception because the coercive pressure itself is
realized adverse treatment.

For positive labels, the model must return verbatim quotations identifying the
victim, the harm, and the corruption--harm connection. The processing script
checks whether the returned evidence occurs in the supplied article text. A
positive prediction supported only by paraphrased, invented, or missing
evidence is mechanically downgraded to \nolinkurl{no_victim}, while the failed
evidence checks remain in the audit output.

For the main models, \nolinkurl{victim_visible} equals 1 for concrete or
institutional/societal victims and 0 for \nolinkurl{no_victim}. Separate
indicators capture concrete and institutional/societal victim visibility.
Articles coded \nolinkurl{unclear} are excluded from analyses using the affected
outcome.

\subsubsection{Corruption Frame}

Corruption frame captures the level at which the article primarily represents
the corruption problem. The labels are \nolinkurl{individualized},
\nolinkurl{systemic}, \nolinkurl{other_or_mixed}, and \nolinkurl{unclear}.
\nolinkurl{individualized} applies when coverage centers on identifiable actors
and a specific allegation, investigation, trial, conviction, resignation, or
scandal. Naming several defendants or briefly mentioning a wider scandal does
not by itself make the frame systemic. \nolinkurl{systemic} applies when the
article's dominant explanation concerns a recurring governance or institutional
pattern, such as state capture, entrenched clientelism, systemic impunity,
rule-of-law breakdown, democratic backsliding, or corruption embedded across
public institutions.

\nolinkurl{other_or_mixed} applies when individualized and systemic frames
receive comparable emphasis, corruption is incidental, or the article is
primarily procedural, administrative, legal, technical, regulatory, or a
roundup of unrelated cases. \nolinkurl{unclear} is reserved for insufficient or
unusably ambiguous text. The coding decision concerns the article's dominant
explanation, not merely whether a person or institution is mentioned. In the
regression models, \nolinkurl{other_or_mixed} is the reference category.

\subsubsection{Case Location}

Case location compares the main corruption case with the publication country
and is coded \nolinkurl{domestic}, \nolinkurl{abroad}, or \nolinkurl{unclear}.
\nolinkurl{domestic} includes national, local, municipal, or sectoral corruption
in the publication country, including domestic misuse of EU funds or domestic
actors' use of offshore accounts. \nolinkurl{abroad} applies when the main actors,
institutions, conduct, or investigation are centered in another country.
Transnational cases are assigned according to their principal center; cases
with several equally central countries and no identifiable center are
\nolinkurl{unclear}. The source agency, quoted speaker, and reporting location do
not determine case location. The derived variable \nolinkurl{abroad_case} equals
1 for \nolinkurl{abroad} and 0 for \nolinkurl{domestic}; unclear cases are excluded
from models using this variable.

\subsubsection{Accused-Actor Visibility}

Accused-actor visibility records whether the article explicitly identifies an
individual or organization as committing, attempting, assisting, enabling,
financing, directing, or concealing the corruption. Allegation or investigation
is sufficient; conviction is not required. Victims, witnesses, whistleblowers,
investigators, prosecutors, judges, regulators, employers, owners, associates,
and beneficiaries do not count unless the article separately alleges their
participation in the corruption.

The labels are \nolinkurl{no_accused_actor}, \nolinkurl{individual_actor},
\texttt{organizational\_or\_}\allowbreak\texttt{institutional\_actor},
\texttt{both\_individual\_and\_}\allowbreak\texttt{organizational}, and
\nolinkurl{unclear}.
Individual and organizational participation are tested independently. An
organization counts only when the article attributes corrupt participation to
the organization itself, for example by stating that it paid bribes,
manipulated a tender, financed a scheme, facilitated money laundering,
concealed misconduct, or was investigated for corruption. An accusation
against an employee, leader, owner, member, subsidiary, or associate does not
automatically implicate the organization. The label
\texttt{both\_individual\_and\_}\allowbreak\texttt{organizational} therefore
requires separate textual support for both actor types.

The derived variable \nolinkurl{accused_actor_visible} equals 1 when at least one
qualifying individual or organizational participant is visible and 0 when no
accused actor is identified. Articles coded \nolinkurl{unclear} are excluded from
models using this variable.

\subsection{Country-Year Context}
\label{sec:method-country-context}

The country-year models include perceived corruption based on Transparency
International's Corruption Perceptions Index (CPI), reverse-coded so that
higher values indicate more perceived corruption:
\[
\texttt{perceived\_corruption}_{ct}
  = 100 - \texttt{CPI}_{ct}.
\]
The main specification uses the lagged value,
$\texttt{perceived\_corruption}_{c,t-1}$, to separate the contextual measure from
current-year coverage. With nine countries, this coefficient is interpreted as
a contextual association rather than a causal country-level effect.
